\documentclass[11pt]{article}
\usepackage[utf8]{inputenc}
\usepackage[margin=1in]{geometry}
\usepackage{amsmath}
\usepackage{amsfonts}
\usepackage{amssymb}
\usepackage{mathtools}
\usepackage{booktabs}
\usepackage{tabularx}
\usepackage{microtype}
\usepackage{needspace}
\usepackage{hyperref}
\usepackage{cleveref}

\setcounter{secnumdepth}{3}
\setcounter{tocdepth}{2}
\setlength{\parskip}{0.5\baselineskip}
\setlength{\parindent}{0pt}
\setlength{\tabcolsep}{5pt}
\newcolumntype{Y}{>{\raggedright\arraybackslash}X}
\widowpenalty=10000
\clubpenalty=10000
\displaywidowpenalty=10000
\allowdisplaybreaks[2]
\setlength{\abovedisplayskip}{8pt plus 2pt minus 2pt}
\setlength{\belowdisplayskip}{8pt plus 2pt minus 2pt}
\setlength{\abovedisplayshortskip}{6pt plus 2pt minus 2pt}
\setlength{\belowdisplayshortskip}{6pt plus 2pt minus 2pt}

\hypersetup{
    colorlinks=true,
    linkcolor=blue,
    citecolor=blue,
    urlcolor=blue
}

\title{Multibody Dynamics: Drift, Control, and the Geometry of Physical Law\\[0.5em]
\large A Unified Tutorial on the Affine Decomposition Across Calculational Frameworks}
\author{}
\date{}

\begin{document}
\maketitle
\tableofcontents
\newpage

%% ============================================================
\section{Introduction: Why Multiple Languages for the Same Physics?}\label{sec:introduction}
%% ============================================================

\subsection{A Motivating Scenario}

Imagine you are designing a controller for a two-jointed robot arm. The arm is in motion---swinging through space with both joints turning. You need to decide, right now, how much torque each motor should apply. But the arm is not sitting still waiting for your command. It is already moving, and its current motion creates forces of its own. Gravity is pulling on it. The rotation of one joint is flinging the other joint around through coupling effects.

This leads to a fundamental question: \textbf{of all the forces and accelerations acting on the system right now, which ones come from the physics (the state the system is already in) and which ones come from the controller (the motor commands)?}

If you can cleanly separate these two contributions, you know what the physics gives you for free---the ``current of the river'' you are floating on---and you know exactly what the motors need to provide.

\subsection{The River Current Analogy}

Think of a kayaker on a river. Even without paddling, the kayaker moves---carried by the river's current. This is \textbf{drift}: the motion the system undergoes on its own, determined entirely by its current state. The paddle strokes are \textbf{control}: deliberate inputs that add to whatever the current is doing. The total motion is always drift plus control.

Now imagine the kayaker asks: ``If I stopped paddling right now, what would happen to me in this instant?'' That question is the \textbf{Zero Torque Counterfactual (ZTCF)}---a snapshot diagnostic, not a prediction of the future.

In mathematical form:
\begin{equation}\label{eq:affine-master}
\dot{x} = f_{\mathrm{drift}}(x) + G(x)\,u,
\end{equation}
where $x$ is the system state, $f_{\mathrm{drift}}(x)$ is the drift (determined by the current state), and $G(x)\,u$ is the control contribution.

\subsection{Three Core Insights}

This article develops three interlocking ideas:

\begin{enumerate}
\item \textbf{Every standard dynamics framework admits a clean drift--control split.} Whether you use energy methods (Euler--Lagrange), force balance on individual bodies (Newton--Euler), spatial screws, task-space coordinates, or abstract geometry, the same structural separation exists.

\item \textbf{Constraint forces are dynamically essential despite doing no net work.} They actively redirect motion and redistribute energy between bodies. They are the guardrails of the dynamical highway---they do not speed you up or slow you down along the road, but they absolutely determine which way the road goes.

\item \textbf{The physics does not depend on the language you use to describe it.} Changing frameworks changes the appearance of equations, not the physical predictions. This enables cross-validation between frameworks.
\end{enumerate}

\subsection{Scope and Assumptions}

This article considers smooth rigid multibody systems with finite-dimensional state descriptions, generalized effort inputs, and standard differential-equation models of motion. Unless otherwise stated, the discussion assumes ideal holonomic or Pfaffian constraints, no impact/contact switching, no actuator internal dynamics, and no dissipative or nonsmooth effects.

\subsection{What Is Standard and What Is New}

The underlying equations are classical. The contribution here is a unified interpretive presentation of the drift/control structure across frameworks, an explicit ZTCF diagnostic as a time-local tool, and a constraint-aware comparison showing how passive and actuated effects appear in each representation.


%% ============================================================
\section{Core Concepts}\label{sec:core-concepts}
%% ============================================================

\subsection{Configuration, State, and Forces}

A system's \textbf{configuration} $q \in \mathbb{R}^n$ describes the spatial arrangement of all parts. The \textbf{state} $x = (q, \dot{q})$ adds velocity information---the minimal data needed to predict future motion.

Generalized effort $\tau$ denotes torques for rotational coordinates and forces for translational coordinates. We distinguish control effort $\tau$, constraint reactions $J_c^T\lambda$, and external loads $\tau_{\mathrm{ext}}$.

\subsection{The Mass Matrix}

The \textbf{mass matrix} $M(q)$ maps applied generalized forces to resulting accelerations. It is symmetric ($M = M^T$) and positive definite, ensuring a unique acceleration for any force.

\textit{Analogy:} Imagine pushing a shopping cart on varying terrain. On tile, a small push gives a big acceleration; on gravel, the same push gives almost nothing. The mass matrix tells you how ``heavy'' the system feels in each direction at each configuration.

\subsection{The Manipulator Equation}

For a rigid multibody system:
\begin{equation}\label{eq:manipulator}
M(q)\,\ddot{q} + C(q,\dot{q})\,\dot{q} + g(q) = \tau,
\end{equation}
where $M(q)\ddot{q}$ is the inertia term (the multibody analogue of $F=ma$), $C(q,\dot{q})\dot{q}$ is the velocity-product term (Coriolis and centripetal effects arising because the mass distribution changes as the system moves), and $g(q) = \partial V/\partial q$ is the gravity term.


%% ============================================================
\section{Running Example: Planar Double Pendulum}\label{sec:double-pendulum}
%% ============================================================

We use a two-link planar pendulum throughout. Let
\begin{equation}\label{eq:dp-coords}
q = \begin{bmatrix}\theta_1 \\ \theta_2\end{bmatrix}, \qquad
\dot{q} = \begin{bmatrix}\dot\theta_1 \\ \dot\theta_2\end{bmatrix},
\end{equation}
with link lengths $(l_1,l_2)$, masses $(m_1,m_2)$, and point masses at each distal end. Angles are measured from the downward vertical ($\theta_1$) and relative to the first link ($\theta_2$).

Endpoint kinematics:
\begin{equation}\label{eq:dp-kinematics}
x_2 = l_1\sin\theta_1 + l_2\sin(\theta_1+\theta_2), \qquad
y_2 = -l_1\cos\theta_1 - l_2\cos(\theta_1+\theta_2).
\end{equation}

The endpoint Jacobian (relating joint velocities to endpoint velocities):
\begin{equation}\label{eq:dp-jacobian}
J(q) = \begin{bmatrix}
l_1\cos\theta_1 + l_2\cos(\theta_1+\theta_2) & l_2\cos(\theta_1+\theta_2) \\
l_1\sin\theta_1 + l_2\sin(\theta_1+\theta_2) & l_2\sin(\theta_1+\theta_2)
\end{bmatrix}.
\end{equation}

The dynamic model is $\tau = M(q)\ddot{q} + C(q,\dot{q})\dot{q} + g(q)$, with:
\begin{equation}\label{eq:dp-mass}
M(q) = \begin{bmatrix}
(m_1+m_2)l_1^2 + m_2l_2^2 + 2m_2l_1l_2\cos\theta_2 & m_2l_2^2+m_2l_1l_2\cos\theta_2 \\
m_2l_2^2+m_2l_1l_2\cos\theta_2 & m_2l_2^2
\end{bmatrix},
\end{equation}
\begin{equation}\label{eq:dp-coriolis}
C(q,\dot{q})\dot{q} = \begin{bmatrix}
-2c\dot\theta_1\dot\theta_2 - c\dot\theta_2^2 \\
c\dot\theta_1^2
\end{bmatrix}, \qquad
c \coloneqq m_2l_1l_2\sin\theta_2,
\end{equation}
\begin{equation}\label{eq:dp-gravity}
g(q) = \begin{bmatrix}
(m_1+m_2)gl_1\sin\theta_1 + m_2gl_2\sin(\theta_1+\theta_2) \\
m_2gl_2\sin(\theta_1+\theta_2)
\end{bmatrix}.
\end{equation}


%% ============================================================
\section{Drift and Control: Separating Physics from Intent}\label{sec:drift-control}
%% ============================================================

\subsection{The Fundamental Decomposition}

Solving the manipulator equation for acceleration:
\begin{equation}\label{eq:el-drift-control}
\ddot{q} = \underbrace{-M^{-1}(q)\bigl[C(q,\dot{q})\dot{q} + g(q)\bigr]}_{\ddot{q}_{\mathrm{drift}}} + \underbrace{M^{-1}(q)\,\tau}_{\ddot{q}_{\mathrm{ctrl}}}.
\end{equation}

The \textbf{drift component} $\ddot{q}_{\mathrm{drift}} = -M^{-1}(C\dot{q} + g)$ is the acceleration with zero actuation: inertial coupling effects ($-M^{-1}C\dot{q}$) plus gravitational loading ($-M^{-1}g$).

The \textbf{control component} $\ddot{q}_{\mathrm{ctrl}} = M^{-1}\tau$ is the incremental acceleration from actuation. The mass matrix $M^{-1}$ acts as the input gain, mapping torques to accelerations.

At fixed state, the decomposition is additive in the input:
\begin{equation}
\ddot{q}_{\mathrm{total}} - \ddot{q}_{\mathrm{drift}} = M^{-1}\tau = \ddot{q}_{\mathrm{ctrl}}.
\end{equation}

This is exact, not approximate. It holds because the input $\tau$ enters linearly.

\subsection{The Zero Torque Counterfactual (ZTCF)}

At time $t$, with state $x(t)$:
\begin{equation}\label{eq:ztcf-def}
\ddot{q}_{\mathrm{ZTCF},t} \coloneqq -M^{-1}\bigl(q(t)\bigr)\bigl[C\bigl(q(t),\dot{q}(t)\bigr)\dot{q}(t) + g\bigl(q(t)\bigr)\bigr].
\end{equation}

This answers: ``If I turned off all motors right now, what acceleration would the system experience at this instant?''

\textbf{ZTCF is time-local.} It is an instantaneous snapshot, not a trajectory prediction. The state evolves and the ZTCF value changes with it.

\textbf{ZTCF reveals accumulated state effects.} A system moving fast has large drift accelerations---not because of current input, but because past inputs built up the current velocity.

\textit{Analogy:} Like a savings account---the current balance (state) is the result of all past deposits and withdrawals. The interest earned this month (drift) depends on the current balance, not on how it was accumulated.

\subsection{State Accumulation}

The current velocity stores the entire past acceleration history:
\begin{equation}\label{eq:state-accumulation}
\dot{q}(t) = \dot{q}(t_0) + \int_{t_0}^{t} \ddot{q}(\sigma)\,d\sigma.
\end{equation}

What appears as ``drift'' at time $t$ is the accumulated result of all prior inputs, gravity, and interactions.

\subsection{Linearization and the Drift--Control Connection}\label{sec:linearization}

Linearizing control-affine dynamics at the ZTCF operating point $(x(t), u = 0)$:
\begin{equation}\label{eq:linearized}
\dot{x} \approx \dot{x}_{\mathrm{ZTCF},t} + A_t\,(x - x(t)) + B_t\,u,
\end{equation}
with $A_t = \partial f_{\mathrm{drift}}/\partial x|_{x(t)}$ and $B_t = G(x(t))$. ZTCF/drift is the zeroth-order anchor; linearization adds first-order sensitivity. This is the foundation for stability analysis, feedback design (LQR, MPC), and observer design.

\subsection{Double-Pendulum Drift--Control Split}

The drift acceleration is:
\begin{equation}
\ddot{q}_{\mathrm{drift}} = -M^{-1}\begin{bmatrix}
-2c\dot\theta_1\dot\theta_2 - c\dot\theta_2^2 + g_1 \\
c\dot\theta_1^2 + g_2
\end{bmatrix},
\end{equation}
with $M^{-1} = (\det M)^{-1}\bigl[\begin{smallmatrix}M_{22}&-M_{12}\\-M_{21}&M_{11}\end{smallmatrix}\bigr]$, $\det M > 0$ always. The control acceleration is $\ddot{q}_{\mathrm{ctrl}} = M^{-1}\tau$, and the total is $\ddot{q} = \ddot{q}_{\mathrm{drift}} + \ddot{q}_{\mathrm{ctrl}}$.

\subsection{Procedure Checklist}

\begin{enumerate}
\item Define generalized coordinates and physical parameters.
\item Derive $T$ and $V$.
\item Extract $M(q)$, $C(q,\dot{q})$, and $g(q)$.
\item Evaluate drift: $\ddot{q}_{\mathrm{drift}} = -M^{-1}(C\dot{q}+g)$ at the current state.
\item Apply control and compute total acceleration.
\item Compare measured acceleration with model prediction.
\item Use residuals to diagnose unmodeled effects.
\end{enumerate}


%% ============================================================
\section{When the System Is Constrained}\label{sec:constrained}
%% ============================================================

\subsection{What Constraints Do}

A holonomic constraint $\phi(q) = 0$ restricts where the system can go. The constraint Jacobian is $J_c = \partial\phi/\partial q$.

\textit{Analogy: guardrails on a highway.} They do zero work along the road (no speeding up or slowing down) but absolutely redirect motion. Constraint forces are the dynamical equivalent of guardrails.

\subsection{Constrained Equations}

\begin{equation}\label{eq:cel-dynamics}
M\ddot{q} + C\dot{q} + g = \tau + J_c^T\lambda + \tau_{\mathrm{ext}},
\end{equation}
\begin{equation}\label{eq:cel-constraint}
J_c\ddot{q} + \dot{J}_c\dot{q} = 0.
\end{equation}

The acceleration-level constraint defines an affine hyperplane of admissible accelerations:
\begin{equation}\label{eq:hyperplane}
\mathcal{H}_t \coloneqq \left\{\ddot{q}\;\middle|\;J_c\ddot{q} = -\dot{J}_c\dot{q}\right\}.
\end{equation}

\subsection{Constrained Drift--Control}

Setting $\tau=0$ gives the constrained drift:
\begin{equation}\label{eq:constrained-drift}
\ddot{q}_{\mathrm{drift}}^{c} = M^{-1}\bigl[-C\dot{q}-g+J_c^T\lambda_0+\tau_{\mathrm{ext}}\bigr],
\end{equation}
where $\lambda_0$ is the constraint multiplier under zero actuation. The constrained control is:
\begin{equation}\label{eq:constrained-ctrl}
\ddot{q}_{\mathrm{ctrl}}^{c} = M^{-1}\bigl[\tau+J_c^T(\lambda-\lambda_0)\bigr].
\end{equation}

The decomposition remains additive: $\ddot{q} = \ddot{q}_{\mathrm{drift}}^{c} + \ddot{q}_{\mathrm{ctrl}}^{c}$.

\subsection{The Tangent-Hyperplane Interpretation}

ZTCF gives a base point on $\mathcal{H}_t$; control moves along allowable directions:
\begin{equation}\label{eq:hyperplane-control}
\ddot{q} = \ddot{q}_{\mathrm{ZTCF},t}^{c} + B_c(x(t))\,\tau,
\end{equation}
where $B_c = M^{-1}[I - J_c^T(J_cM^{-1}J_c^T)^{-1}J_cM^{-1}]$ is the constraint-consistent input map. This is exact within the stated assumptions.

\subsection{Why Constraint Forces Matter Despite Zero Work}

Ideal constraint reactions do not contribute net generalized power on admissible motions, but they remain dynamically essential. They enforce compatibility, redirect motion, and transmit internal loads through which momentum and energy are redistributed across the coupled system. In whip-like motions, these reactions are the mechanism by which kinetic energy migrates from proximal to distal segments.

\subsection{Procedure Checklist}

\begin{enumerate}
\item Define $\phi(q)=0$.
\item Compute $J_c$ and $\dot{J}_c\dot{q}$.
\item Build $\mathcal{H}_t$.
\item Solve the KKT system for $(\ddot{q},\lambda)$.
\item Compute $\ddot{q}_{\mathrm{ZTCF},t}^{c}$ by setting $\tau=0$.
\item Recover $J_c^T\lambda$; decompose into passive and actuation-induced parts.
\item Compare constrained and unconstrained drift baselines.
\item Audit per-coordinate power flow.
\end{enumerate}


%% ============================================================
\section{The Same Physics in Different Languages}\label{sec:other-frameworks}
%% ============================================================

Having established the drift--control decomposition in the Euler--Lagrange framework, we now show how the same ideas appear in four other formulations. Each is like a different map of the same city, emphasizing different features.

\subsection{Newton--Euler: Body-Level Force Analysis}\label{sec:newton-euler}

Newton--Euler is the formulation that most directly exposes how accelerations are supported by local force and moment transmission through the chain. For each body $i$:
\begin{equation}\label{eq:ne-trans}
m_i a_{c,i} = f_i^{\mathrm{in}} - f_i^{\mathrm{out}} + f_i^{\mathrm{ext}} + f_i^{\mathrm{ctrl}},
\end{equation}
\begin{equation}\label{eq:ne-rot}
I_i\dot\omega_i + \omega_i\times(I_i\omega_i) = n_i^{\mathrm{in}} - n_i^{\mathrm{out}} + n_i^{\mathrm{ext}} + n_i^{\mathrm{ctrl}}.
\end{equation}

The recursive algorithm (forward pass: base to tip; backward pass: tip to base) makes internal load transfer explicit. Joint reaction forces are visible here but hidden inside the mass matrix coupling in Euler--Lagrange.

ZTCF: set $\tau=0$ and solve the recursion for drift accelerations and passive joint reactions.

\subsection{Product-of-Exponentials and Screw Theory}\label{sec:poe-screw}

Screw theory provides a unified language for spatial velocity and force by combining rotational and translational quantities into six-dimensional objects: \textbf{twists} (velocity) and \textbf{wrenches} (force).

Forward kinematics on $SE(3)$:
\begin{equation}\label{eq:poe-fk}
g_{sb}(q) = e^{\hat{S}_1 q_1}e^{\hat{S}_2 q_2}\cdots e^{\hat{S}_n q_n}g_{sb}(0).
\end{equation}

Spatial twist: $\mathcal{V}_s = J_s(q)\dot{q}$. Spatial acceleration: $\dot{\mathcal{V}}_s = J_s\ddot{q} + \dot{J}_s\dot{q}$, where $\dot{J}_s\dot{q}$ is the bias acceleration (screw-coordinate analogue of Coriolis/centripetal effects).

The velocity-dependent term $\mathrm{ad}_{\mathcal{V}}^{*}(\mathcal{I}\mathcal{V})$ expresses the same physics as Coriolis/centripetal terms in the language of Lie algebras.

For the double pendulum, screw axes $S_1 = [0,0,1,0,0,0]^T$ and $S_2 = [0,0,1,0,-l_1,0]^T$. The space Jacobian recovers the same endpoint velocity as~\eqref{eq:dp-jacobian}.

\subsection{Operational-Space Analysis}\label{sec:operational-space}

Operational-space analysis answers: ``What does the system look like from the task's perspective?'' Define task coordinates $x_{\mathrm{task}}=\psi(q)$ with Jacobian $J=\partial\psi/\partial q$.

Task-space dynamics:
\begin{equation}\label{eq:os-dynamics}
\Lambda(q)\ddot{x}_{\mathrm{task}}+\mu(q,\dot{q})+p(q) = F_{\mathrm{ctrl}}+F_{\mathrm{constr}}+F_{\mathrm{ext}},
\end{equation}
where $\Lambda(q)=(JM^{-1}J^T)^{-1}$ is the task-space inertia (valid where $JM^{-1}J^T$ is nonsingular), $\mu$ is the task-space Coriolis/centripetal term, and $p$ is the task-space gravity.

The task-space drift is the image of joint-space drift under the kinematic map:
\begin{equation}\label{eq:os-drift}
\ddot{x}_{\mathrm{drift}} = J\ddot{q}_{\mathrm{drift}} + \dot{J}\dot{q}.
\end{equation}

The relation $\tau = J^T F$ expresses virtual-work equivalence. The reverse association from joint torques to task-level effects is generally non-unique in redundant systems.

Singularities occur when the Jacobian loses rank ($\det J = 0$); near these configurations, the task-space inertia becomes singular.

\subsection{Coordinate-Free Geometric Mechanics}\label{sec:geometric}

The geometric formulation works directly with intrinsic objects on the configuration manifold $Q$ and its tangent bundle $TQ$, making representation independence visible.

\textit{Analogy:} The Earth is round, and no flat map perfectly represents it. Geometric mechanics works directly with the ``globe'' rather than any particular flat map.

On $Q$ with kinetic-energy Riemannian metric:
\begin{equation}\label{eq:geo-eom}
\nabla_{\dot{q}}\dot{q} + \mathrm{grad}\,V(q) = \mathcal{G}(q)u + \mathcal{F}_{\mathrm{ext}} + \mathcal{F}_{\mathrm{constr}}.
\end{equation}

The covariant acceleration $\nabla_{\dot{q}}\dot{q}$ expands to $\ddot{q}^i + \Gamma^i_{jk}\dot{q}^j\dot{q}^k$ in coordinates, where Christoffel symbols $\Gamma^i_{jk}$ are exactly the Coriolis/centripetal coefficients. These arise from the Levi-Civita connection, not directly from curvature; nonzero Christoffel symbols can exist in flat spaces (e.g., polar coordinates in Euclidean space).

Intrinsic drift on $TQ$:
\begin{equation}\label{eq:geo-drift}
f_{\mathrm{drift}}(q,\dot{q}) = \begin{pmatrix}\dot{q}\\-\Gamma(\dot{q},\dot{q})-\mathrm{grad}\,V(q)+\mathcal{F}_{\mathrm{ext}}\end{pmatrix}.
\end{equation}

For the double pendulum, $Q = S^1\times S^1$ (a torus). The drift field matches Euler--Lagrange exactly. Coordinate independence is verified by switching charts: different coordinate expressions, same physical accelerations.


%% ============================================================
\section{Cross-Formalism Equivalence}\label{sec:cross-formalism}
%% ============================================================

\subsection{Equivalent Representations of the Same Balance Laws}

For the same mechanism at the same state:
\begin{align}
\text{Euler--Lagrange:}\quad & M\ddot{q}+C\dot{q}+g-J_c^T\lambda = \tau+\tau_{\mathrm{ext}}, \\
\text{Screw:}\quad & \mathcal{I}\dot{\mathcal{V}}+\mathrm{ad}_{\mathcal{V}}^*(\mathcal{I}\mathcal{V})+\mathcal{W}_g-\mathcal{W}_{\mathrm{constr}} = \mathcal{W}_{\mathrm{ctrl}}+\mathcal{W}_{\mathrm{ext}}, \\
\text{Op-space:}\quad & \Lambda\ddot{x}+\mu+p = F_{\mathrm{ctrl}}+F_{\mathrm{constr}}+F_{\mathrm{ext}}, \\
\text{Geometric:}\quad & \nabla_{\dot{q}}\dot{q}+\mathrm{grad}\,V = \mathcal{G}u+\mathcal{F}_{\mathrm{ext}}+\mathcal{F}_{\mathrm{constr}}.
\end{align}

Equivalence means that under appropriate kinematic, coordinate, and frame mappings, the formulations produce the same physical motion, the same power balance, and compatible descriptions of reactions.

\subsection{ZTCF Anchor in Each Formalism}

At fixed state:
\begin{itemize}
\item Euler--Lagrange: evaluate $\ddot{q}$ with $\tau=0$.
\item Newton--Euler: run recursion with control removed.
\item Screw: evaluate spatial dynamics with zero control wrenches.
\item Operational-space: map passive acceleration to task space.
\item Geometric: evaluate intrinsic drift field.
\end{itemize}

\subsection{Frame Changes and Invariants}

Twists transform as $V_s = \mathrm{Ad}_g\, V_b$, wrenches as $F_s = \mathrm{Ad}_g^{-T} F_b$. Instantaneous power is invariant: $F_s^T V_s = F_b^T V_b$.


%% ============================================================
\section{What the Equations Hide: Loads and Energy Transfer}\label{sec:hidden-loads}
%% ============================================================

Compact equations hide joint reactions, internal energy transfer, and contact force distributions. Recovery strategy: introduce explicit channels and solve for reaction variables.

\subsection{Power-Channel Decomposition}

\begin{equation}
P_{\mathrm{ctrl}}=\tau^T\dot{q}, \qquad
P_{\mathrm{drift}}=(C\dot{q}+g)^T\dot{q}, \qquad
P_{\mathrm{constr}}=(J_c^T\lambda)^T\dot{q}=0.
\end{equation}

Even when total constraint power is zero, per-coordinate power $P_{\mathrm{constr},i}=[J_c^T\lambda]_i\dot{q}_i$ can be large---only the sum vanishes. Per-coordinate power decompositions are useful bookkeeping devices but depend on the choice of coordinates; only total power balance is invariant.

In whip-like motions, constraint forces extract energy from proximal links and inject it into distal links, producing speed amplification.


%% ============================================================
\section{Symmetry, Invariance, and Conservation Laws}\label{sec:symmetry}
%% ============================================================

The equations of motion are objective: they do not depend on coordinates.

\textit{Analogy:} The distance between two cities does not depend on whether you measure it in miles or kilometers. The physical distance is real; the unit system is a convention.

By Noether's theorem:
\begin{itemize}
\item Translation symmetry $\Rightarrow$ linear momentum conservation.
\item Rotation symmetry $\Rightarrow$ angular momentum conservation.
\item Time-translation symmetry $\Rightarrow$ energy conservation.
\end{itemize}

Conservation laws interact with the drift--control decomposition: if a quantity is conserved under drift, then only control can change it. If energy is conserved under drift, the drift dynamics preserve energy surfaces and control moves between them.

These are structural properties of the physics, not of one coordinate system or framework. They provide framework-independent cross-validation checks.


%% ============================================================
\section{Practical Guidance}\label{sec:practical}
%% ============================================================

\begin{center}
\small
\renewcommand{\arraystretch}{1.2}
\begin{tabularx}{\textwidth}{>{\raggedright\arraybackslash}p{3.1cm}Y Y}
\toprule
Formalism & Best question answered & When to use \\
\midrule
Euler--Lagrange & How inputs change generalized accelerations & Energy-based analysis, control design \\
\midrule
Constrained E--L & How constraints redirect motion and loads & Reaction force analysis \\
\midrule
Newton--Euler & Where local loads and transfer pathways occur & Structural analysis, load paths \\
\midrule
PoE / screw & How motion and force compose across frames & Spatial analysis, frame transformations \\
\midrule
Operational-space & How task objectives map to joint effort & End-effector control \\
\midrule
Coordinate-free & Which conclusions are representation-invariant & Theoretical verification \\
\bottomrule
\end{tabularx}
\end{center}

Common pitfalls: confusing drift with static equilibrium, assuming constraint forces are negligible, treating ZTCF as a trajectory prediction, forgetting that operational-space quantities depend on full joint configuration, and assuming per-coordinate power decompositions are frame-invariant.


%% ============================================================
\section{Conclusion}\label{sec:conclusion}
%% ============================================================

Under the assumptions stated in \cref{sec:introduction}, each formulation considered in this article admits a natural separation between passive state-determined dynamics and input-dependent contributions.

Ideal constraint reactions do not contribute net generalized power on admissible motions, but they remain dynamically essential because they enforce compatibility, redirect motion, and transmit internal loads through which momentum and energy are redistributed across the coupled system.

The formulations considered here are best viewed as equivalent representations of the same underlying balance laws, rather than as literally identical equations written in different notation.

Used carefully, ZTCF provides a compact time-local reference for distinguishing what the current state would produce on its own from what the actuators are contributing at that instant.

\end{document}
